\documentclass[english,utf8,aspectratio=169]{beamer}
\usetheme{tud}
\usecolortheme{tud}
\usepackage[backend = biber,style = numeric]{biblatex}
\addbibresource{bibliography.bib}
\usepackage[inkscapepath={./build/svg-inkscape}]{svg}
\svgpath{{./figures/}}
\usepackage{hyperref}
\usepackage{booktabs}
\usepackage{algorithm}
\usepackage[noEnd=false]{algpseudocodex}
\usepackage{pgfplots}
\pgfplotsset{compat=1.18}
\usepackage{pgfplotstable}
\usepackage{pgffor}

\AtBeginSection[]{\partpage{\usebeamertemplate***{part page}}}
\setbeamertemplate{theorems}[numbered]

\title{Vinculum: A Reusable Implementation of Lock-Free Deduplication}
\subtitle{Bachelor's Thesis Defense}
\author{Eric Niklas Wolf}
\institut{Chair of Distributed and Networked Systems}
\date{April 1, 2025}

\begin{document}
    \maketitle

    \begin{frame}
        \frametitle{Outline}
        \tableofcontents
    \end{frame}

    \section{Motivation}
    \begin{frame}
        \frametitle{Motivation}
        \begin{itemize}
            \item periodic backups are important for protecting against data loss
            \item most backups share large amounts of data with each other
            \item when eliminating redundant data it is desirable to do so across
                the greatest number of backups possible
            \item existing solutions are either proprietary or have limitations
        \end{itemize}

        \begin{figure}
            \begin{tabular}{ c c c c }
                \toprule
                & Original size & Compressed size & Deduplicated size \\
                \midrule
                Daily backup & 460.65 GB & 405.49 GB & 454.43 MB \\
                \midrule
                All backups & 19.58 TB & 16.73 TB & 745.75 GB \\
                \bottomrule
            \end{tabular}
            \caption{\small{Example of \href{https://www.borgbackup.org}{BorgBackup} removing redundant data.}}
        \end{figure}
    \end{frame}

    \section{Background}
    \begin{frame}
        \frametitle{Terminology}
        \begin{itemize}
            \item Client: Machine creating backups
            \item Repository: Storage for backups which is shared by multiple clients
            \item Chunk: Part of a backup stored in the repository
            \item Manifest: Metadata associated with a specific backup which references chunks
        \end{itemize}

        \begin{center}
            \includesvg[height=.5\textheight, pretex=\small]{terminology.drawio.svg}
        \end{center}
    \end{frame}
    
    \begin{frame}
        \frametitle{Lock-Free Deduplication}
        \begin{itemize}
            \item after deleting manifests unreferenced chunks should be pruned
            \item doing so while a manifest is being created can cause data loss
            \item this is prevented by performing the deletion in two steps
        \end{itemize}

        \begin{center}
            \includesvg[width=\textwidth, pretex=\small]{lock-free_deduplication.drawio.svg}
        \end{center}
    \end{frame}

    \begin{frame}
        \frametitle{Example}
        \begin{center}
            \includesvg[height=.7\textheight, pretex=\small]{lock-free_deduplication_example.drawio.svg}
        \end{center}
    \end{frame}

    \section{Implementation}
    \begin{frame}
        \frametitle{Goals}
        \begin{itemize}
            \item assist authors of backup programs in implementing Lock-Free Deduplication

                $\rightarrow$ handles non-atomic operations, error recovery and client management

            \item provide enough flexibility to allow for different solutions to use cases
                like encryption, compression, \ldots

                $\rightarrow$ focuses on the algorithm only
            
            \item provide common interfaces on which these other use cases can be implemented

                $\rightarrow$ dropped because designing good universal interfaces is difficult
        \end{itemize}
    \end{frame}

    \begin{frame}
        \frametitle{Vinculum}
        \begin{itemize}
            \item \textit{Repository and Manifest} traits: describe operations
                used by the Lock-Free Deduplication algorithm
            \item \textit{FossilCollection} struct: stores data required by the
                fossil deletion step and provides method to perform it
            \item \textit{FossilCollectionBuilder} struct: provides methods to perform
                the fossil collection step
            \item \textit{PipelinedFossilCollectionBuilder} struct: allows reuse data
                aquired during a fossil deletion step for a fossil collection step
        \end{itemize}
    \end{frame}

    \begin{frame}
        \frametitle{Fossil collection step Pseudocode}
        {
            \small
            \begin{algorithmic}[1]
                \State $builder \gets \Call{createBuilder}$
                \ForAll{$manifest \in \Call{existingManifestFiles}{repository}$}
                    \If{$manifest \notin manifestFilesToDelete$}
                        \State $chunks \gets \Call{downloadChunkReferences}{manifest}$
                        \State \Call{markAsReferenced}{chunks, builder}
                        \State \Call{markAsSeen}{manifest, builder}
                    \EndIf
                \EndFor
                \ForAll{$manifest \in manifestFilesToDelete$}
                    \State $chunks \gets \Call{downloadChunkReferences}{manifest}$
                    \State \Call{markAsFossilCandidates}{chunks, builder}
                \EndFor
                \State $fossilCollection \gets \Call{collectFossils}{builder, repository}$
                \State \Call{deleteManifestFiles}{manifestFilesToDelete}
            \end{algorithmic}
        }
    \end{frame}

    \begin{frame}
        \frametitle{Pipelining}
        \begin{center}
            \includesvg[width=\textwidth, pretex=\small]{pipelining.drawio.svg}
        \end{center}
    \end{frame}

    \section{Performance Evaluation}
    \begin{frame}
        \frametitle{Setup}
        \begin{itemize}
            \item a simple repository was implemented ontop a UNIX file system
                together with a command line application
            \item three existing backup programs with a similar internal architekture
                where selected for comparison 
            \item the evaluation was performed on a machine running Linux 6.12
                with the ext4 file system located on a hard disk
        \end{itemize}
        \begin{center}
            \begin{tabular}{ c c c c c }
                \toprule
                \textbf{Program} & Borg 1 & Borg 2 & Duplicacy \\
                \midrule
                \textbf{Version} & 1.4.0  & 2.0.0b14 & 3.2.4 \\
                \midrule
                \textbf{Programming language} & Python, C & Python, C & Go \\
                \bottomrule
            \end{tabular}
        \end{center}
    \end{frame}

    \begin{frame}
        \frametitle{VM Images: Setup}
        \begin{itemize}
            \item three VM images where backed up with a fixed chunk size of 1 MiB
            \begin{enumerate}
                \item after completing the setup
                \item after installing common developer tools
                \item after a power on immediately followed by a power off
            \end{enumerate}
            \item they where deleted in FIFO order
            \item all programs performed the deletion with two operations which where measured independently
        \end{itemize}
    \end{frame}

    \begin{frame}
        \frametitle{VM Images: Results}
        \pgfplotstableread[col sep = comma]{benchmark/results/vm_images.csv}{\vmimagedata}
        \begin{tikzpicture}
            \pgfplotsset{every axis/.style={
                ybar stacked,
                ymin=0,ymax=170,
                xtick=data,
                xtick style={draw=none},
                symbolic x coords={
                    Borg 1,
                    Borg 2,
                    Duplicacy,
                    Vinculum
                },
                bar width=8pt,
                width=.9\textwidth,
                height=.95\textheight
            }}

            \begin{axis}[hide axis,bar shift=-10pt]
                \addplot table [x=tool,y=ubuntu_setup1] {\vmimagedata};
                \addplot table [x=tool,y=ubuntu_setup2] {\vmimagedata};
            \end{axis}

            \begin{axis}[hide axis]
                \addplot table [x=tool,y=ubuntu_tools1] {\vmimagedata};
                \addplot table [x=tool,y=ubuntu_tools2] {\vmimagedata};
            \end{axis}

            \begin{axis}[
                bar shift=10pt,
                ylabel=Runtime in seconds,
                name=mainplot
            ]
                \addplot table [x=tool,y=ubuntu_onoff1] {\vmimagedata};
                \addplot table [x=tool,y=ubuntu_onoff2] {\vmimagedata};
            \end{axis}

            % subplot
            \begin{axis}[
                hide axis,
                bar shift=-10pt,
                restrict x to domain=1:3,
                at={(mainplot.north east)},
                anchor=north east,
                xshift=-.5cm,
                yshift=-.5cm,
                ymax=3,
                small
            ]
                \addplot table [x=tool,y=ubuntu_setup1] {\vmimagedata};
                \addplot table [x=tool,y=ubuntu_setup2] {\vmimagedata};
            \end{axis}

            \begin{axis}[
                hide axis,
                restrict x to domain=1:3,
                at={(mainplot.north east)},
                anchor=north east,
                xshift=-.5cm,
                yshift=-.5cm,
                ymax=3,
                small
            ]
                \addplot table [x=tool,y=ubuntu_tools1] {\vmimagedata};
                \addplot table [x=tool,y=ubuntu_tools2] {\vmimagedata};
            \end{axis}

            \begin{axis}[
                bar shift=10pt,
                restrict x to domain=1:3,
                at={(mainplot.north east)},
                anchor=north east,
                xshift=-.5cm,
                yshift=-.5cm,
                ymax=3,
                minor y tick num=1,
                small
            ]
                \addplot table [x=tool,y=ubuntu_onoff1] {\vmimagedata};
                \addplot table [x=tool,y=ubuntu_onoff2] {\vmimagedata};
            \end{axis}
        \end{tikzpicture}
    \end{frame}

    \begin{frame}
        \frametitle{Chunk Scaling: Setup}
        \begin{itemize}
            \item a specific number of chunks is divided between four backups
            \item one backup is deleted afterwards
            \item the sum of the times required by the two operations required for deletion was measured
        \end{itemize}
    \end{frame}

    \foreach \f in {4096, 32768}{
        \begin{frame}
            \frametitle{Chunk Scaling: Results}
            \pgfplotstableread[col sep = comma]{benchmark/results/chunks\f .csv}{\chunkdata}
            \begin{tikzpicture}
                \pgfplotsset{
                    width=.9\textwidth,
                    height=.85\textheight
                }
                \begin{axis}[
                    xlabel=Number of chunks (\f{} bytes),
                    ylabel=Runtime in seconds,
                    width=0.75\textwidth,
                    legend pos=outer north east,
                    scaled x ticks=base 10:-3
                ]
                    \addplot table [x=chunks, y expr=\thisrow{Borg 1 deletion}+\thisrow{Borg 1 compaction}] {\chunkdata};
                    \addlegendentry{Borg 1}
                    \addplot table [x=chunks, y expr=\thisrow{Borg 2 deletion}+\thisrow{Borg 2 compaction}] {\chunkdata};
                    \addlegendentry{Borg 2}
                    \addplot table [x=chunks, y expr=\thisrow{Duplicacy collection}+\thisrow{Duplicacy deletion}] {\chunkdata};
                    \addlegendentry{Duplicacy}
                    \addplot table [x=chunks, y expr=\thisrow{Vinculum collection}+\thisrow{Vinculum deletion}] {\chunkdata};
                    \addlegendentry{Vinculum}
                \end{axis}
            \end{tikzpicture}
        \end{frame}
    }

    \begin{frame}
        \frametitle{Backups Scaling: Setup}
        \begin{itemize}
            \item a fixed number of chunks with a size of 4096 bytes where created
            \item they are divided between a specific number of backups
            \item one of the backups is deleted afterwards
            \item the sum of the times required by the two operations required for deletion was measured
        \end{itemize}
    \end{frame}

    \foreach \f in {0, 1000}{
        \begin{frame}
            \frametitle{Backups Scaling: Results}
            \pgfplotstableread[col sep = comma]{benchmark/results/manifests\f const.csv}{\manifestdata}
            \begin{tikzpicture}
                \pgfplotsset{
                    width=.9\textwidth,
                    height=.85\textheight
                }
                \begin{axis}[
                    xlabel=Number of backups (\f{} chunks),
                    ylabel=Runtime in seconds,
                    width=0.75\textwidth,
                    legend pos=outer north east,
                    name=mainplot
                ]
                    \addplot table [x=manifests, y expr=\thisrow{Borg 1 deletion}+\thisrow{Borg 1 compaction}] {\manifestdata};
                    \addlegendentry{Borg 1}
                    \addplot table [x=manifests, y expr=\thisrow{Borg 2 deletion}+\thisrow{Borg 2 compaction}] {\manifestdata};
                    \addlegendentry{Borg 2}
                    \addplot table [x=manifests, y expr=\thisrow{Duplicacy collection}+\thisrow{Duplicacy deletion}] {\manifestdata};
                    \addlegendentry{Duplicacy}
                    \addplot table [x=manifests, y expr=\thisrow{Vinculum collection}+\thisrow{Vinculum deletion}] {\manifestdata};
                    \addlegendentry{Vinculum}
                \end{axis}

                %subplot
                \begin{axis}[
                    cycle list shift=2,
                    at={(mainplot.north west)},
                    anchor=north west,
                    xshift=.75cm,
                    yshift=-.5cm,
                    scaled x ticks=base 10:-2,
                    tiny
                ]
                    \addplot table [x=manifests, y expr=\thisrow{Duplicacy collection}+\thisrow{Duplicacy deletion}] {\manifestdata};
                    \addplot table [x=manifests, y expr=\thisrow{Vinculum collection}+\thisrow{Vinculum deletion}] {\manifestdata};
                \end{axis}
            \end{tikzpicture}
        \end{frame}
    }

    \section{Discussion}
    \begin{frame}
        \frametitle{Realism of the Performance Evaluation}
        \begin{itemize}
            \item command line application representing Vinculum is very simple
                (only fixed size chunking, no compression, encryption or caching)
            \item despite choosing similar settings the performance of the other
                programs may be influenced by implementing these features
            \item comparing them under real world scenarios would require additional
                features to be implemented
        \end{itemize}
    \end{frame}

    \begin{frame}
        \frametitle{Missing Features}
        \begin{itemize}
            \item variable-size chunking: allows deduplication when inserting or removing bytes
            \item compression: saved around 10\% of space in motivational example
            \item encryption: preserves privacy when using cloud storage
            \item caching: important when performing more expensive chunking operations
                or network accesses
        \end{itemize}

        \begin{figure}
            \begin{tabular}{ c c }
                \toprule
                Caches & Duration \\
                \midrule
                chunks \& files & 0.8s \\
                \midrule
                chunks & 16.3s \\
                \bottomrule
            \end{tabular}
            \caption{\small{Example of \href{https://www.borgbackup.org}{BorgBackup} benefitting from caching.}}
        \end{figure}
    \end{frame}

    \section{Conclusion and Future Work}
    \begin{frame}
        \frametitle{Conclusion}
        \begin{itemize}
            \item a reusable implementation of Lock-Free Deduplication was presented
            \item it performed favorably in benchmarks when compared with existing solutions
            \item there is still much to do to create a fully functional backup program
        \end{itemize}
    \end{frame}

    \begin{frame}
        \frametitle{Future Work}
        \begin{itemize}
            \item implement missing features (preferably as independent components)
            \item investigate robustness against the repository experiencing power failure
            \item investigate alternative consistency models for possible performance improvements
            \item review security\footnote{\url{https://www.daemonology.net/blog/chunking-attacks.pdf}}
                of variable-size chunking
        \end{itemize}
    \end{frame}
    
    \section{Questions}
\end{document}
